\documentclass[12pt]{article}
%\renewcommand\normalsize{\fontsize{10pt}{12pt}\selectfont}
\usepackage{graphicx} % Required for inserting images
\usepackage[left=2cm,right=2cm,top=2cm,bottom=1.5cm]{geometry}

\title{HW 3}
%\author{Dennis Wilkinson}
%\date{March 31 2025; due Monday April 7}

\begin{document}

\setlength{\parindent}{0pt}
\setlength{\parskip}{3mm}

FS1 Wilkinson HW 3
Circular motion

1. Rope swing. I linked to a random rope swing video in the email that announced this HW, featuring someone named Adam at Barton Creek. Adam has to grab the rope tightly, because it can be hard to hang on as one swings through the bottom; it definitely requires more force than just hanging motionless.

a. Hanging motionless: Adam's arms have to support his weight. Try to guess his mass, and then calculate the force he would need to exert to hang motionless.

b. Swinging: Adam is now moving in a circular path; for this he requires a centripetal acceleration. Estimate Adam's velocity at the bottom of the swing, and the rope length, to calculate his centripetal acceleration. Ballpark numbers fine here.

c. His arms must be providing the force for the centripetal acceleration. Calculate this using your numbers from (a) and (b). Reminder: he must still be balancing his weight force; the centripetal force is in addition to this. Compare the force here to that of part (a); what's the percentage increase? 

d. How does the swinging speed and rope length impact how hard it is to hang on? Explain.


2. Let's say Earth, mass $m$, is in a circular orbit around the sun, mass $M$. (It's not quite circular, but that's okay.)

a. Calculate Earth's velocity with respect to the center of the sun, assuming the sun is motionless during the orbit. The angular velocity $\omega$ will be $2\pi$ divided by the length of a year, in seconds. You can look up $r$; this length is called an AU but you need it in meters.

b. The centripetal force for Earth is provided by gravity between Earth and the Sun, $F=GMm/r^2$. Show that as a result, $v=\sqrt{\frac{GM}{r}}$. Use your number from (a) to calculate $G$ and see how close you get to the real value. Look up $M$ for the sun.

c. Neptune has a very nearly circular orbit with radius $\sim$30.1 AU (Earth's radius is 1 AU) and mass $\sim$17.15 times $m_{Earth}$. How does its velocity compare to Earth's? Also calculate its year, in Earth years. 

3* (non-calculus, but a bit tricky). A roller coaster goes around a loop-the-loop of radius $r$. Calculate the minimum velocity required at the top of the loop so that the cars stay in contact with the track and don't fall down. Hint: there is a contact force between the car and the track, which can help provide the centripetal acceleration. 

4. Read  the article (link in the email) about Kepler, Brahe, and the ``war on Mars."

5. Briefly summarize the article from \#4; two or three paragraphs ought to do it.
\end{document}